\documentclass[10pt,a4paper]{amsart}
\usepackage[margin=19mm]{geometry}
\usepackage{amsmath,amssymb,mathtools}
\usepackage[scr=rsfs,cal=euler]{mathalfa}
\usepackage{xfrac}
\usepackage[unicode,hidelinks,bookmarks=false]{hyperref}
\newcommand{\ie}{i.e.\ }
\newcommand{\cf}{cf.\ }
\newcommand{\resp}{respectively\ }
\newcommand{\Subsection}{\subsection}
\providecommand{\nobreakdash}{\nobreak\hbox{-}}
\setlength{\emergencystretch}{2em}
\multlinegap0pt
\usepackage{bm}
\usepackage{amssymb}
\usepackage[shortlabels]{enumitem}
\usepackage{tikz}
\usepackage[normalem]{ulem}
\newtheorem{lemma}{Lemma}
\usepackage{cleveref}
\crefname{lemma}{Lemma}{Lemmas}
\renewcommand{\Pr}{\mathop{\mathrm{Pr}}\nolimits}
\newcommand{\defeq}{:=}
\newcommand{\tp}[1]{{\left( #1 \right)}}
\title{Rapid mixing in positively weighted restricted Boltzmann machines}
\author{Author(s), Weiming Feng, Heng Guo, Minji Yang}
\date{Source: arXiv:2604.00963v2 (CC BY 4.0)}
\begin{document}
\maketitle

\noindent\emph{Source and licence.} Rapid mixing in positively weighted restricted Boltzmann machines. Version arXiv:2604.00963v2 (CC BY 4.0), distributed by arXiv under the Creative Commons Attribution 4.0 International licence, http://creativecommons.org/licenses/by/4.0/, as recorded in the arXiv licence metadata for this exact version and shown on the abstract page https://arxiv.org/abs/2604.00963v2. Full licence notice: \texttt{licenses/arxiv-2604.00963-CC-BY-4.0.txt}.\par\smallskip

\noindent\textit{Editorial scope.} Counted unit: the article's lemma lemma:warm-start-configuration-general (source0.tex lines 2862--2870). The article's complete original proof of it is delivered with it (source0.tex lines 2878--3023, one \texttt{proof} environment of 146 lines, which the article places away from the statement). No mathematics is written, repaired or re-derived: the statement and the proof are the article's own lines, in the article's own notation, and no retained line is substituted or re-flowed.  Retained uncounted as the source-local context the proof needs: lines 700--702.  Nothing was omitted from any retained span, so the package carries no omission record.  No retained line contains a citation, so the wrapper carries no bibliography and no bibliography entry.  No retained line contains a figure environment, an image-inclusion command or drawing code, so no image is delivered and the package declares no asset dependency.  Editorial formatting only, in addition: an AMS \texttt{amsart} preamble that restates the article's own declarations and macros used by the retained lines, namely the environments \texttt{lemma} on separate counters, together with the standard packages those lines need.  The retained environments print in this document as Lemma 1 in order of appearance, whereas the article prints the same items with the numbering of its own full document.  The complete native source of the article as distributed in the arXiv e-print for this exact version is bundled unchanged as \texttt{sources/197623/source0.tex}.
\begin{align}\label{eq:mixing-time-glauber-gap}
t_{\textnormal{mix}}^P(\epsilon) \leq \frac{1}{\gamma_{\text{GD}}}\log \frac{1}{\epsilon^2 \mu(\sigma)}.
\end{align}

\begin{lemma}\label{lemma:warm-start-configuration-general}
Let $\lambda > 0$ be a constant. 
Let $\mu$ be a Gibbs distribution of a ferromagnetic two-spin system on a graph $G = (V,E)$ with $(\beta_e,\gamma_e)_{e\in E}$ and $(\lambda_v)_{v \in V}$. 
Suppose $\beta_e \gamma_e > 1$ and $\beta_e \leq 1 \leq \gamma_e$ for all $e \in E$, and $\lambda_v < \lambda$ for all $v \in V$.
Let $P_{\mu}^{\textnormal{Glauber}}$ be the Glauber dynamics on $\mu$. Suppose the spectral gap of the Glauber dynamics is at least $0 < g < 1$.  Then, the mixing time of the Glauber dynamics on $\mu$ satisfies
\begin{align*}
t_{\textnormal{mix}}^{\textnormal{Glauber}}\tp{\frac{1}{4e}} \leq O_{\lambda}\left(|V| \log |V| + \frac{|V|^2}{g} \log |V|\right).
\end{align*}
\end{lemma}

\begin{proof}[Proof of \Cref{lemma:warm-start-configuration-general}]
Let $N = |V|$. 
Let $N_0(\lambda)$ be a sufficiently large constant depending only on $\lambda$. First we consider the case when $N \leq N_0(\lambda) =  O_\lambda(1)$.
In each update of the Glauber dynamics, we have a chance at least $\frac{1}{1+\lambda}$ to update the value of a vertex to 1. 
We run the Glauber dynamics for some $O_\lambda(1)$ steps, so that with probability $\Omega_\lambda(1)$, all vertices take the value 1. 
Let $T_0 = O_\lambda(1)$ be a sufficiently large constant. With probability at least $1 - \frac{1}{10e}$, we can find a time $t < T_0$ such that all vertices take the value 1.
For each edge, $\gamma_e > 1 \geq \beta_e$. It holds that $\mu(\boldsymbol{1}) = \Omega_\lambda(1)$ if $N \leq N_0(\lambda)$. Using~\eqref{eq:mixing-time-glauber-gap}, starting from all-1 configuration, we only need to run Glauber dynamics for $O_\lambda(1/g)$ steps to get a configuration with total variation distance at most $\frac{1}{10e}$ to the stationary distribution $\mu$. 
A simple coupling argument shows that the total variation distance between the resulting distribution and $\mu$ is at most $\frac{1}{4e}$ after $T_0 + O_\lambda(1/g) = O_\lambda(1/g)$ steps.

Now, we assume $N \geq N_0(\lambda)$ is large enough.
Fix $\tau \in \{0,1\}^V$.
We say that a vertex $u \in V$ is \emph{bad} in $\tau$ if
\begin{align*}
    \lambda_u \leq \frac{1}{100 N^5}
    \qquad\text{and}\qquad
    \tau(u) = 0.
\end{align*}
For any edge $e = \{u,w\} \in E$, we say that $e$ is \emph{bad} in $\tau$ if
\begin{align*}
\gamma_e \geq 100N^5 \text{ and } (\tau(u) = 0 \text{ or } \tau(w) = 0),
\end{align*}
and we say that $\tau$ is a \emph{warm-start configuration} if no vertex or edge is bad in $\tau$.

We prove the following two claims.
\begin{itemize}
    \item Starting from an arbitrary configuration $X_0 \in \{0,1\}^{V}$, after running $P_{\mu}^{\textnormal{Glauber}}$ for $T_0 = O_\lambda(N (\log N)^2)$ steps, with probability at least $1 - \frac{1}{10e}$, the configuration $X_{T_0}$ is a warm-start configuration.
    \item Starting from any warm-start configuration $X_{T_0}$, after running the Glauber dynamics for $T_1 = O_\lambda\left(\frac{N^2}{g} \log N\right)$ steps, where $g$ is a lower bound of the spectral gap, the total variation distance between the resulting distribution and $\mu$ is at most $\frac{1}{10e}$.
\end{itemize}
If these two claims hold, we can construct a coupling between the law of $X_{T_0 + T_1}$ and the stationary distribution $\mu$ such that the coupling fails with probability at most
\begin{align*}
&\Pr[X_{T_0} \text{ is not a warm-start configuration}] + \Pr[\text{coupling fails} \mid X_{T_0} \text{ is a warm-start configuration}] \\
&= \frac{1}{10e} + \frac{1}{10e} < \frac{1}{4e},
\end{align*}
which finishes the proof.

Now we prove the first claim. Let $M = C_1 N \log N$ and $L = C_0 \log N$, where $C_1 > 0$ is a sufficiently large absolute constant and $C_0 = C_0(\lambda) > 0$ is a sufficiently large constant depending only on $\lambda$. Set
\begin{align*}
T_0 = LM = O_\lambda(N(\log N)^2).
\end{align*}
Partition the time interval $[T_0]$ into $L$ consecutive blocks, each of length $M$.
We list the sequence of updated vertices as
\begin{align*}
  v_1,v_2,\ldots,v_{T_0}.
\end{align*}
An update sequence is \emph{good} if every vertex is updated at least once in every block.
By the coupon collector bound and a union bound over all $L$ blocks, the update sequence is good with probability at least $1 - \frac{1}{20e}$.

Fix a good update sequence.
We first bound the probability that a vertex is bad in $X_{T_0}$. Fix any vertex $u \in V$, and let $t_u$ be the last time at which $u$ is updated. We must have $\lambda_u \leq \frac{1}{100N^5}$, since otherwise $u$ cannot be bad. Fix all the updates before time $t_u$. Let $u_1,u_2,\ldots,u_d$ denote all neighbors of $u$, let $\beta_i,\gamma_i$ denote the parameters of the edge $\{u_i,u\}$, and let $\rho$ denote the configuration of the other variables at time $t_u$. Then
\begin{align}\label{eq:lambda-u-bound}
 \frac{\Pr[ u \text{ is updated to } 0]}{\Pr[u \text{ is updated to }1]} = \lambda_u \prod_{i \in [d]: \rho(u_i) = 1} \frac{1}{\gamma_i}\prod_{i \in [d]: \rho(u_i) = 0} \beta_i.
\end{align}
Since $\frac{1}{\gamma_i} \leq 1$ and $\beta_i \leq 1$, we have
\begin{align*}
 \Pr[X_{t_u}(u) = 0] \leq \lambda_u \leq \frac{1}{100N^5}.
\end{align*}
Hence, the probability that $u$ is bad in $X_{T_0}$ is at most \smash{$\frac{1}{100N^5}$}.

Now fix a bad edge $e = \{u,w\} \in E$ with  $\gamma_e \geq 100N^5$, as otherwise, $e$ cannot be bad.
We call a pair of times $(t,t')$ a \emph{clean pair} for $e$ if $t < t'$, $\{v_t,v_{t'}\} = \{u,w\}$, $v_t \neq v_{t'}$, and for all $t < \ell < t'$ we have $v_\ell \notin \{u,w\}$.
Since the update sequence is good, both $u$ and $w$ are updated at least once in every block. Fix any block. %Restrict the update sequence in this block to the times at which either $u$ or $w$ is updated. 
Since both vertices appear in the block, there must be a clean pair of times for $e$. 
We list all clean pairs in the update sequence: $(t_j,t'_j)_{j=1}^K$ with $t'_j < t_{j+1}$,
where $K \geq L \geq C_0 \log N$ since there is at least one clean pair in each block.

Fix all randomness used to update vertices in $V \setminus \{u,w\}$.
Let $p_\lambda \defeq \frac{1}{1+\lambda}$. For each clean pair $(t_b,t_b')$, define the event $A_b$ by
\begin{align*}
A_b = \{X_{t_b'}(u) = X_{t_b'}(w) = 1\}.
\end{align*}
By~\eqref{eq:lambda-u-bound}, at every update of either $u$ or $w$, the chosen vertex is updated to $1$ with probability at least $p_\lambda$. Therefore, conditional on all past updates on $\{u,w\}$ before time $t_b$, the probability of the event $A_b$ is at least $p_\lambda^2$.
By iterating this bound over all clean pairs and choosing $C_0$ sufficiently large as a function of $\lambda$, we obtain
\begin{align*}
\Pr\Big[\bigcap_{b = 1}^K \overline{A_b}\Big] \leq (1-p_\lambda^2)^K \leq N^{-6}.
\end{align*}
If any event $A_b$ occurs, then the following event $A$ holds:
\begin{itemize}
\item $A$: there exists the first time $t_e < T_0$ such that $X_{t_e}(u) = X_{t_e}(w) = 1$.
\end{itemize}
Hence, $t_e$ exists with probability at least $1 - N^{-6}$. Furthermore, the random variable $t_e$ % is fixed by all updates before or at the time $t_e$ and
is independent from the updates after $t_e$.
Suppose $t_e = s'$ and let $s > s'$ be the first time after $t_e$ at which the edge $u$ or $w$ is updated to $0$.
At time $s$, one of the endpoints, say $u$, is updated to $0$ while the other endpoint is still equal to $1$. Hence, by~\eqref{eq:lambda-u-bound},
\begin{align*}
\Pr[X_s(u) = 0 \mid X_{s-1}(w) = 1] \leq \frac{\lambda_u}{\gamma_e} \leq \frac{\lambda}{100N^5}.
\end{align*}
A union bound over all times $s'$ for $t_e = s'$ and all times $s$ for $s > s'$ yields
\begin{align*}
\Pr[e \text{ is bad in } X_{T_0} \mid A] \leq \frac{\lambda T_0^2}{100N^5}.
\end{align*}
Therefore, since $T_0 = O_\lambda(N(\log N)^2)$, we have
\begin{align*}
\Pr[e \text{ is bad in } X_{T_0}] \leq \Pr[\neg A] + \Pr[e \text{ is bad in } X_{T_0} \mid A] \leq N^{-6} + \frac{\lambda T_0^2}{100N^5} \leq \frac{1}{100N^{2.5}}.
\end{align*}

Taking a union bound over all vertices and edges, conditioned on the update sequence fixed above, the probability that $X_{T_0}$ is not a warm-start configuration is at most
\begin{align}\label{eq:warm-start-configuration-probability}
\frac{N}{100N^5} + \frac{N^2}{100N^{2.5}} < \frac{1}{20e}.
\end{align}
Combining this with the probability $\frac{1}{20e}$ that the update sequence is not good proves the first claim.

For the second claim, we show a lower bound on $\mu(\tau)$ for each warm-start configuration $\tau$. For any configuration $\tau' \in \{0,1\}^{V}$, not necessarily a warm-start configuration, we give a lower bound on the ratio $\frac{\mu(\tau)}{\mu(\tau')}$. We analyze the contribution of every vertex and every edge in $G$. Formally, the ratio $\frac{\mu(\tau)}{\mu(\tau')}$ can be written as the following ratio of products:
\begin{align*}
  \frac{\mu(\tau)}{\mu(\tau')} = \frac{ \prod_{u \in V} a_u(\tau(u)) \prod_{e \in E} b_e(\tau(e)) }{ \prod_{u \in V} a_u(\tau'(u)) \prod_{e \in E} b_e(\tau'(e)) },
\end{align*}
where, for each vertex $u \in V$,
\begin{align*}
  a_u(\tau(u))\defeq
  \begin{cases}
    \lambda_u & \text{ if }\tau(u) = 0;\\
    1 & \text{ if }\tau(u) = 1,
  \end{cases}
\end{align*}
and for each edge $e = \{u,w\} \in E$,
\begin{align*}
  b_e(\tau(e))\defeq
  \begin{cases}
    \beta_e & \text{ if }\tau(u) = \tau(w) = 0;\\
    \gamma_e & \text{ if }\tau(u) = \tau(w) = 1;\\    
    1 & \text{ if }\tau(u) \neq \tau(w).
  \end{cases}
\end{align*}
We analyse each ratio as follows.
\begin{itemize}
    \item If $\lambda_u \leq \frac{1}{100N^5}$, then $\tau(u) = 1$ because $\tau$ is warm-start. Hence, $\frac{a_u(\tau(u))}{a_u(\tau'(u))} \geq \min\{1,\lambda^{-1}\}$.
    \item If $\lambda_u > \frac{1}{100N^5}$, then $\frac{a_u(\tau(u))}{a_u(\tau'(u))} \geq \min\{1/(100N^5),\lambda^{-1}\}$.
    \item If $\gamma_e \geq 100N^5$, then $\tau(u) = \tau(w) = 1$ because $\tau$ is warm-start. Therefore, $\frac{b_e(\tau(e))}{b_e(\tau'(e))} \geq 1$.
    \item If $\gamma_e < 100N^5$, then $\beta_e > \frac{1}{\gamma_e} > \frac{1}{100N^5}$ because $\beta_e \gamma_e > 1$. Therefore,
    \begin{align*}
        \frac{b_e(\tau(e))}{b_e(\tau'(e))} \geq \frac{\beta_e}{\gamma_e} > \frac{1}{10^4 N^{10}}.
    \end{align*}
\end{itemize}
The total number of edges in $E$ is at most $N^2$. Hence, the ratio $\frac{\mu(\tau)}{\mu(\tau')}$ can be bounded as follows:
\begin{align*}
\frac{\mu(\tau)}{\mu(\tau')} \geq \tp{\min\{1/(100N^5),\lambda^{-1}\}}^{N} \cdot \tp{\frac{1}{10^4 N^{10}}}^{N^2} \geq \exp(-O_\lambda(N^2 \log N)).
\end{align*}
Since the above lower bound holds for every $\tau' \in \{0,1\}^V$, summing over all $2^N$ choices of $\tau'$ gives
\begin{align}\label{eq:mu-tau-lower-bound}
 \mu(\tau) \geq \exp(-O_\lambda(N^2 \log N)) \cdot 2^{-N} = \exp(-O_\lambda(N^2 \log N)).
\end{align}
Let
\begin{align*}
 T_1 \defeq O\left( \max_{\text{warm-start } \tau} \frac{1}{g} \log \frac{1}{(1/10e)^2 \mu(\tau)}\right) = O_\lambda\tp{\frac{N^2}{g} \log N}.
\end{align*}
The second claim follows from~\eqref{eq:mixing-time-glauber-gap} with $\epsilon = \frac{1}{10e}$ for the warm-start configuration $\tau$.
\end{proof}

\end{document}
